\section{Short-orbit CRT and ambient integer incidence}
\label{sec:tpc35-integer-incidence}

Large ultra divisors give strong uniqueness along a \(J\)-long orbit.
This section separates that fact from a different question: how many
distinct divisor atoms can land at the same orbit time or integer point.

\subsection{An exact short-orbit incidence norm}

Let \(I\subset\Z\) be an interval of diameter strictly smaller than every
modulus \(u_x\) in a finite atom set \(\cX\).  Associate to atom \(x\) one
residue class \(r_x\pmod{u_x}\), and define the incidence matrix
\begin{equation}
 B_{x,j}=\one_{j\equiv r_x\,({\rm mod}\ u_x)},
 \qquad x\in\cX,\quad j\in I.
 \label{eq:tpc35-short-orbit-matrix}
\end{equation}

\begin{proposition}[Exact short-orbit norm]
\label{prop:tpc35-short-orbit-norm}
Every row of \(B\) has at most one nonzero entry, and
\begin{equation}
 \boxed{
 \|B\|_{\op}^2
 =\max_{j\in I}\#\{x\in\cX:j\equiv r_x\pmod{u_x}\}.}
 \label{eq:tpc35-short-orbit-norm}
\end{equation}
\end{proposition}

\begin{proof}
Two distinct integers in \(I\) differ by less than \(u_x\), so they cannot
belong to the same residue class modulo \(u_x\).  Thus every row contains
at most one \(1\).  It follows that \(B^*B\) is diagonal, and its \(j\)-th
diagonal entry is the number of atoms incident to \(j\).  Its largest
eigenvalue is the right side of \eqref{eq:tpc35-short-orbit-norm}.
\end{proof}

For a physical atom \((m,u)\), primitivity gives
\((u,hm)=1\) and the residue is \(r\equiv-hm^{-1}\pmod u\).  Since
\(u>T>J\) at the endpoint, the proposition applies to every ultra atom.
For two atoms, \cref{prop:tpc35-affine-CRT-compatibility} determines
whether their residue classes have a joint return modulo the least common
multiple.  The remaining difficulty is precisely the column degree in
\eqref{eq:tpc35-short-orbit-norm}: uniqueness for each atom does not bound
how many different atoms choose the same time.

\subsection{An integer fixed-determinant transform}

Let \(h\ne0\) be fixed.  Consider primitive coefficient pairs
\(
 \lambda=(a,s)
\)
in a box where the \(a\)-interval has length at most \(A\) and
\(1\le s\le S\).  Associate the integer line
\begin{equation}
 L_{a,s}=\{(x,y)\in\Z^2:sy-ax=h\},
 \qquad (a,s)=1.
 \label{eq:tpc35-integer-line}
\end{equation}
Let \(\cP\subset\Z_{>0}^2\) be contained in a box whose \(x\)-interval
has length at most \(X_0\) and whose points satisfy \(y>T\).  Define
\begin{equation}
 (\cR F)(a,s)=\sum_{P\in L_{a,s}\cap\cP}F(P).
 \label{eq:tpc35-integer-Radon}
\end{equation}
For a point \(P\), let
\begin{equation}
 r(P)=\#\{(a,s):P\in L_{a,s}\}.
 \label{eq:tpc35-point-degree}
\end{equation}

\begin{lemma}[Primitive-line Schur bound]
\label{lem:tpc35-line-Schur}
One has
\begin{equation}
 \boxed{
 \|\cR\|_{\op}^2
 \le\max_{(a,s)}\sum_{P\in L_{a,s}\cap\cP}r(P).}
 \label{eq:tpc35-line-Schur}
\end{equation}
\end{lemma}

\begin{proof}
The kernel of \(\cR\cR^*\) is the nonnegative intersection count
\[
 (\cR\cR^*)((a,s),(a',s'))
 =|L_{a,s}\cap L_{a',s'}\cap\cP|.
\]
Its row sum is
\(
 \sum_{P\in L_{a,s}\cap\cP}r(P)
\).  Schur's test proves the result.
\end{proof}

Two distinct primitive lines have at most one rational intersection.  If
\(a's-as'\ne0\), it is
\begin{equation}
 x=\frac{h(s'-s)}{a's-as'},
 \qquad
 y=\frac{h(a'-a)}{a's-as'}.
 \label{eq:tpc35-line-intersection}
\end{equation}
The next theorem gives the corresponding integral degree bound without
discarding parallel or nonintegral cases.

\begin{theorem}[Ambient integer-incidence norm]
\label{thm:tpc35-integer-incidence}
With the preceding hypotheses, every point \(P=(x,y)\in\cP\) satisfies
\begin{align}
 r(x,y)
 &\ll_h1+\min\left\{\frac Ay,\frac Sx\right\}
 \le_h1+\frac AT,
 \label{eq:tpc35-point-degree-bound}\\
\intertext{and every line \(L_{a,s}\) in the coefficient family satisfies}
 |L_{a,s}\cap\cP|
 &\le1+\frac{X_0}{s}.
 \label{eq:tpc35-line-point-bound}
\end{align}
Consequently
\begin{equation}
 \boxed{
 \|\cR\|_{\op}^2
 \ll_h\left(1+\frac{X_0}{s_{\min}}\right)
 \left(1+\frac AT\right),}
 \label{eq:tpc35-integer-incidence-bound}
\end{equation}
where \(s_{\min}\) is the least \(s\) in the coefficient family.
\end{theorem}

\begin{proof}
Fix \(P=(x,y)\).  Any two coefficient solutions of
\(sy-ax=h\) differ by
\begin{equation}
 \Delta s=\frac{x}{(x,y)}k,
 \qquad
 \Delta a=\frac{y}{(x,y)}k,
 \qquad k\in\Z.
 \label{eq:tpc35-coefficient-solution-step}
\end{equation}
If \(P\) lies on one of the lines, then \((x,y)\mid h\), so
\((x,y)\le|h|\).  Counting the admissible \(k\)'s in the \(a\)- and
\(s\)-intervals gives the first estimate in
\eqref{eq:tpc35-point-degree-bound}; \(y>T\) gives the second.

Now fix a primitive line.  Its integer points satisfy
\[
 ax\equiv-h\pmod s.
\]
Because \((a,s)=1\), this is one residue class for \(x\pmod s\).  An
interval of length \(X_0\) contains at most \(1+X_0/s\) such integers,
and \(y\) is then unique.  This proves \eqref{eq:tpc35-line-point-bound}.
Insert the two bounds into \cref{lem:tpc35-line-Schur} and maximize over
the line family.
\end{proof}

\subsection{Physical containing-box comparison}

On a projected affine layer from TPC-33, the coefficient and point boxes
have scales
\begin{equation}
 A\asymp LJv,\qquad X_0\asymp D/v,
 \qquad 1\le v\le D,
 \label{eq:tpc35-physical-box-scales}
\end{equation}
and the reflected complementary slope has \(s\ge1\).  Hence
\begin{align}
 \|\cR\|_{\op}^2
 &\ll_h
 \left(1+\frac{D}{s_{\min}v}\right)
 \left(1+\frac{LJv}{T}\right)
 \notag\\
 &\ll_h1+\frac XT
 \label{eq:tpc35-physical-incidence-scale}
\end{align}
at the endpoint.  Indeed \(LJD=QJ\asymp X\), the cross term is at most
\(X/T\), and \(D\le X/T\) there.

This containing-box bound is sharp in scale.  Take \(v\asymp D\), a point
\(x=1,y\asymp T\), and coefficients
\begin{equation}
 a=sy-h.
 \label{eq:tpc35-pencil-example}
\end{equation}
For \(\gg_h X/T\) values of \(s\le cX/T\) coprime to \(h\), these are
primitive lines in an \(a\)-box of length \(O(X)\) passing through the
same point.  Thus one can have \(r(P)\gg_h X/T\) in the ambient box.

After extracting the outer row \(L^2\) mass \(Q\), the uniform orbit-Gram
ledger would require an unsigned incidence constant of scale \(Q/J\).
The ratio supplied by the containing-box relaxation is
\begin{equation}
 \frac{X/T}{Q/J}=\frac{J^2}{T}
 =X^{279/1000+o(1)},
 \label{eq:tpc35-incidence-gap}
\end{equation}
the same gap as \eqref{eq:tpc35-shared-ultra-gap}.  This does not prove
that the physical incidence transform is sharp: the actual point set is
selected by squarefree and M\"obius weights, gcd projectors, smooth
supports, and coupled row labels.  It says that the unsigned ambient box
alone reproduces, rather than removes, the shared-ultra loss.
No norm comparison between the actual masked orbit Gram matrix and this
containing-box transform is asserted.
